\section{Delete Operations for Two Strings}
\label{sec:dp-delete-operations}

\problemstatement{Given two strings $s_1$ (length $n$) and $s_2$ (length
$m$), find the minimum number of character \textbf{deletions} (from
either string) needed to make the two strings \textbf{equal}.}

\begin{patternbox}
This is Section~\ref{sec:dp-edit-distance}'s Edit Distance restricted to
\textbf{only the delete operation} -- and that restriction collapses it
into a direct reduction to Section~\ref{sec:dp-lcs}'s LCS, rather than
needing its own three-way recurrence at all.
\end{patternbox}

\subsection*{Understanding the reduction}

If the only allowed operation is deletion, the two strings can only ever
become equal by each reducing down to some common subsequence of both --
deleting everything that isn't part of that shared subsequence. To
minimize total deletions, we want to keep as much as possible, which
means keeping the \textbf{longest} common subsequence and deleting
everything else:
\[
  \textit{minDeletions} = (n - |\textit{LCS}(s_1,s_2)|) + (m -
  |\textit{LCS}(s_1,s_2)|).
\]
The first term is how many characters of $s_1$ are deleted (everything
not part of the kept LCS); the second is the same for $s_2$.

\subsection*{All Four Approaches}

Every stage reuses Section~\ref{sec:dp-lcs}'s LCS implementation
unchanged; only this final arithmetic combination is new.

\begin{lstlisting}
#include <bits/stdc++.h>
using namespace std;

int lcsSpaceOptimized(string& s1, string& s2) {
    int n = s1.size(), m = s2.size();
    vector<int> prev(m + 1, 0), current(m + 1, 0);

    for (int i = 1; i <= n; i++) {
        for (int j = 1; j <= m; j++) {
            if (s1[i - 1] == s2[j - 1]) {
                current[j] = 1 + prev[j - 1];
            } else {
                current[j] = max(prev[j], current[j - 1]);
            }
        }
        prev = current;
    }
    return prev[m];
}

int minDeletions(string& s1, string& s2) {
    int lcsLen = lcsSpaceOptimized(s1, s2);
    return (s1.size() - lcsLen) + (s2.size() - lcsLen);
}

int main() {
    string s1 = "sea", s2 = "eat";
    cout << minDeletions(s1, s2) << endl; // 2
    return 0;
}
\end{lstlisting}

\subsection*{Dry run}

\dryrun{Take $s_1=\texttt{"sea"}$, $s_2=\texttt{"eat"}$.}

LCS of \texttt{"sea"} and \texttt{"eat"} is \texttt{"ea"} (length $2$).
$\textit{minDeletions} = (3-2)+(3-2) = 1+1 = 2$: delete \texttt{'s'}
from $s_1$ (leaving \texttt{"ea"}) and delete \texttt{'t'} from $s_2$
(leaving \texttt{"ea"}) -- both strings become \texttt{"ea"} with $2$
total deletions.

\begin{complexitybox}
\textbf{Time Complexity.} $O(n \cdot m)$, inherited directly from LCS.
\textbf{Space Complexity.} $O(m)$ for the rolling-row LCS.
\end{complexitybox}

\begin{takeawaybox}
Whenever an edit-distance-style problem restricts the allowed operations
to a strict subset (here, deletions only), check whether the restricted
version reduces to a simpler, already-solved problem rather than
re-deriving a new three-way (or however-many-way) recurrence. ``Delete
only, from both strings, to reach equality'' is exactly ``maximize the
shared kept subsequence,'' i.e.\ LCS, in disguise.
\end{takeawaybox}
